\section{Qualitative estimates for the finite Gram test}
\label{app:quantitative}

This appendix supplies the qualitative estimates used in
\Cref{lem:negative-limit}.  The main theorem only requires their
large-$m$ limits.  We also record, at the end of the appendix, the
quantitative constants used to obtain the explicit sufficient order in
\Cref{app:explicit-bound}.

\subsection{The scaled limit of the Gram-invariant part}

Let
\[
 \delta_m:=m^{-1/2},\qquad
 \mathsf x_\ell^{(m)}
 =(\xi_\ell^{(m)},\eta_\ell^{(m)})
 =\left(ie^{-\delta_m a_\ell},
        -ie^{-\delta_m\overline{a_\ell}}\right),
 \qquad 1\le\ell\le s.
\]
The speed parameters $a_1,\ldots,a_s$ are fixed throughout, with
$\Re a_\ell>0$.  Let $\mathbf L_{\infty,m}$ denote the matrix obtained
from $\mathbf L_m$ by extending the depth indices in $D,K,B$ from
$\{0,\ldots,m-1\}$ to all nonnegative integers.

For the pair $(\ell,h)$, set
\[
\begin{aligned}
 A_m^{\ell h}
 &:=(1-\xi_\ell^{(m)}\overline{\xi_h^{(m)}})
    (1-\eta_\ell^{(m)}\overline{\eta_h^{(m)}}),\\
 C_m^{\ell h}
 &:=(1-\xi_\ell^{(m)}\eta_\ell^{(m)})
    (1-\overline{\xi_h^{(m)}}\,\overline{\eta_h^{(m)}}),\\
 \Delta_m^{\ell h}
 &:=(1-\xi_\ell^{(m)}\overline{\eta_h^{(m)}})
    (1-\eta_\ell^{(m)}\overline{\xi_h^{(m)}}).
\end{aligned}
\]
Direct summation of the geometric series associated with $D,K,B$ gives
\begin{equation}\label{B:eq:Linfty}
 (\mathbf L_{\infty,m})_{\ell h}
 =
 2\left(
 \frac{2}{(A_m^{\ell h})^2}
 +\frac{1}{(C_m^{\ell h})^2}
 -\frac{1}{A_m^{\ell h}C_m^{\ell h}}
 -\frac{2}{(\Delta_m^{\ell h})^2}
 \right).
\end{equation}

\begin{lemma}[Finite-depth tail]\label{B:lem:tail}
For the fixed finite set of test points, there exist constants
$C,c>0$, independent of $m$, such that
\[
 \max_{1\le\ell,h\le s}
 \left|
 (\mathbf L_m-\mathbf L_{\infty,m})_{\ell h}
 \right|
 \le C m^4 e^{-c\sqrt m}.
\]
Consequently,
\[
        \frac{1}{m^2}
        \bigl(\mathbf L_m-\mathbf L_{\infty,m}\bigr)
        \longrightarrow0.
\]
\end{lemma}

\begin{proof}
All coefficients in the finite sums defining $D,K,B$ and their
realignments grow at most polynomially, of degree four, in the depth
indices.  Since the set of speeds is fixed and
$\rho_0:=\min_\ell\Re a_\ell>0$, every test-point monomial of total
depth $q$ has modulus at most $e^{-\rho_0 q/\sqrt m}$.  Every term
discarded when the infinite sums are truncated at depth $m$ contains
at least one index of size at least $m$.  Summing the resulting
polynomially weighted geometric tails gives
\[
 C m^4e^{-c\sqrt m}
\]
for suitable constants $C,c>0$.  Division by $m^2$ proves the second
assertion.
\end{proof}

For the fixed speed parameters define
\begin{equation}\label{B:eq:AC-limit}
 A_{\ell h}:=
 |a_\ell+\overline{a_h}|^2,\qquad
 C_{\ell h}:=
 4\,\Re(a_\ell)\Re(a_h).
\end{equation}
Since $1-e^{-\delta z}=\delta z+O(\delta^2)$ as $\delta\to0$,
uniformly for the present finite set of $z$'s,
\[
 m A_m^{\ell h}\longrightarrow A_{\ell h},\qquad
 m C_m^{\ell h}\longrightarrow C_{\ell h},\qquad
 \Delta_m^{\ell h}\longrightarrow4.
\]
Combining this with \eqref{B:eq:Linfty} and
\Cref{B:lem:tail} gives
\[
 \frac{1}{m^2}(\mathbf L_m)_{\ell h}
 \longrightarrow
 2\left(
 \frac{2}{A_{\ell h}^2}
 +\frac{1}{C_{\ell h}^2}
 -\frac{1}{A_{\ell h}C_{\ell h}}
 \right).
\]
Since $s$ is fixed, the convergence is also convergence in every
matrix norm.  This is the limit matrix $\mathbf L_\star$ used in
\Cref{lem:negative-limit}.

\subsection{Qualitative control of the Gram-dependent part}

We now prove the only uniform estimate needed in the main theorem:
for the fixed finite test,
\[
 \sup_{\widehat Q_{\rm av}\succeq0}
 \frac{1}{m^2}
 \max_{\ell,h}|(\mathbf R_m)_{\ell h}|
 \longrightarrow0.
\]
The argument has two ingredients.  Positivity first eliminates all
low-total mixed modes.  This gives an improved estimate on conjugate
points.  A qualitative maximum-principle argument then propagates the
improvement to the finitely many realigned points. 
For $0\le g\le2m-2$, let
\[
 s_g:=\sum_{\substack{p+q=g\\0\le p,q<m}}\mathbf e_{pq}
\]
in the mixed coordinates.

\begin{lemma}[Forced low mixed modes]\label{B:lem:means}
The baseline mixed block $B$ is positive semidefinite, and every
positive semidefinite normal-form Gram representation satisfies
\begin{equation}\label{B:eq:means}
                 Ms_g=Bs_g=0,\qquad 0\le g<m.
\end{equation}
\end{lemma}

\begin{proof}
The matrix $B$ is block diagonal with respect to the total
$g=p+q$.  On a complete block, its off-diagonal entries are
nonpositive and its row sums vanish; hence the block is a weighted
graph Laplacian and is positive semidefinite.  Incomplete blocks are
principal submatrices, so $B\succeq0$.

To obtain the second assertion, substitute the corner variables so that
the pure wedges vanish and the mixed wedge $w_{pq}$ equals $t^{p+q}$.
If $Q\succeq0$ is any feasible Gram matrix, the resulting polynomial
is
\[
 \left\|\sum_g t^gQ^{1/2}s_g\right\|^2.
\]
The same polynomial computed from the baseline mixed block has no term
of degree below $2m$.  If $g<m$ were the least index for which
$Q^{1/2}s_g\ne0$, the coefficient of $t^{2g}$ would be strictly
positive, a contradiction.  Thus $Qs_g=0$ for $g<m$.  Restricting to
the mixed block gives \eqref{B:eq:means}.
\end{proof}

Write
\[
        \Psi_m:=M-B.
\]
By \Cref{B:lem:means}, the repeated generating expression associated
with $\Psi_m$ contains only total degrees at least $m$ in each pair of
variables.  At the conjugate points used below this high-order support
is exponentially suppressed.  Together with the explicit baseline
block $D-K$, this yields the following qualitative estimate.

\begin{lemma}[Conjugate-point improvement]\label{B:lem:conjugate}
Let $\alpha,\beta$ range over fixed compact real intervals with
$\alpha>0$, and set
\[
 \xi_m=ie^{-(\alpha+i\beta)/\sqrt m},\qquad
 \eta_m=-ie^{-(\alpha-i\beta)/\sqrt m}.
\]
Uniformly over all positive semidefinite normal-form Gram
representations,
\[
        P_-[\xi_m,\eta_m]=O(m),
\]
whereas on every fixed complex neighbourhood on which the two node
moduli stay strictly inside the unit disk at distance comparable with
$m^{-1/2}$,
\[
        P_-[\xi,\eta]=O(m^2).
\]
\end{lemma}

\begin{proof}
The global estimate follows from positivity and the diagonal pure
block:
\[
 0\le P_-[\xi,\eta]
 \le
 \frac{8}{(1-|\xi|^2)^2(1-|\eta|^2)^2}.
\]
At the scaled points above, the right-hand side is $O(m^2)$.

On the conjugate slice, the same-side four-form vanishes and the
Pl\"ucker correction can be rewritten through $\Psi_m$.  The baseline
term $D-K$ is $O(m)$ there, while the high-total support of $\Psi_m$
from \Cref{B:lem:means} contributes $o(m)$.  Hence
$P_-[\xi_m,\eta_m]=O(m)$ uniformly.  A direct tail estimate proving the last $o(m)$ statement is obtained exactly as in \Cref{B:lem:tail}; the quantitative version is recorded below.
\end{proof}

We use the following standard qualitative form of the two-constant
principle.

\begin{lemma}[Qualitative continuation]\label{B:lem:continuation}
Let $W_m$ be a holomorphic vector-valued function on a fixed rectangle
$\Omega\subset\C$, possibly with a range dimension depending on $m$.
Assume $\|W_m\|\le C$ on $\Omega$ and
$\|W_m\|\to0$ uniformly on a nontrivial open segment $I$ of one side
of the rectangle.  Then $\|W_m\|\to0$ uniformly on every compact subset
of $\Omega$ having positive harmonic measure with respect to $I$.
\end{lemma}

\begin{proof}
For nonzero $W_m$, the function $\log\|W_m\|$ is subharmonic; zeros are
handled by the usual value $-\infty$.  The two-constant theorem bounds
$\log\|W_m(z)\|$ by the boundary majorant weighted by the harmonic
measure of $I$.  On a fixed compact subset this harmonic measure is
bounded below by a positive constant.  Since the bound on $I$ tends to
zero while the global bound remains fixed, the resulting interior bound
also tends to zero.  The identically zero case is immediate.
\end{proof}

\begin{proposition}[Uniform vanishing of the Gram-dependent term]
\label{B:prop:R}
For the fixed test points of \eqref{eq:test-nodes},
\begin{equation}\label{B:eq:Rvanish}
 \sup_{\widehat Q_{\rm av}\succeq0}
 \frac1{m^2}\max_{1\le\ell,h\le s}
 |(\mathbf R_m)_{\ell h}|
 \longrightarrow0.
\end{equation}
Consequently, for every fixed $c\in\C^s$,
\[
 \sup_{\widehat Q_{\rm av}\succeq0}
 \left|
 c^*\frac{\mathbf R_m}{m^2}c
 \right|
 \longrightarrow0.
\]
\end{proposition}

\begin{proof}
Let $P_-=\Gamma^\top\Gamma$ and define
\[
 \mathbf h_m(\alpha,\beta)
 :=
 \frac1m\,
 \Gamma f\left(
 ie^{-(\alpha+i\beta)/\sqrt m},
 -ie^{-(\alpha-i\beta)/\sqrt m}
 \right).
\]
By \Cref{B:lem:conjugate}, $\|\mathbf h_m\|$ is uniformly bounded on
fixed complex rectangles, while it tends uniformly to zero on the real
conjugate slice.

Apply \Cref{B:lem:continuation} first in the $\alpha$ variable and then
in the $\beta$ variable.  No explicit harmonic-measure constant is
needed because the test set is finite and fixed.  For each pair
$(\ell,h)$, put
\[
 \alpha_{\ell h}:=
 \frac{a_\ell+\overline{a_h}}2,\qquad
 \beta_{\ell h}:=
 \frac{a_\ell-\overline{a_h}}{2i}.
\]
These finitely many points lie in a fixed compact set.  Realignment
gives
\[
 \frac1{m^2}(\mathbf R_m)_{\ell h}
 =
 \left\|
 \mathbf h_m(\alpha_{\ell h},\beta_{\ell h})
 \right\|^2,
\]
and therefore every entry tends to zero, uniformly over all feasible
Gram representations.  Since there are only finitely many entries,
\eqref{B:eq:Rvanish} follows.  Finally,
\[
 \left|
 c^*\frac{\mathbf R_m}{m^2}c
 \right|
 \le
 \|c\|_1^2\,
 \frac1{m^2}\max_{\ell,h}|(\mathbf R_m)_{\ell h}|,
\]
which proves the last assertion.
\end{proof}

\subsection{Explicit estimates for the sufficient threshold}
\label{B:subsec:effective}

For completeness, we now record the quantitative version of the preceding
argument that is used in \Cref{app:explicit-bound}.  These constants are not
needed for the qualitative proof in the main text.

For four arguments in the open unit disk define
\begin{equation}\label{B:eq:denominators}
\begin{gathered}
 A=(1-\xi\zeta)(1-\eta\omega),\qquad
 \Delta=(1-\xi\omega)(1-\eta\zeta),\\
 C=(1-\xi\eta)(1-\zeta\omega),\qquad
 \Pi=\xi\eta\zeta\omega.
\end{gathered}
\end{equation}
At the level of scalar generating expressions, write
$Q_{0,m}:=D-K$ and let $L_m$ denote the Gram-invariant kernel
corresponding to \(\mathbf L_m\).  A subscript \(\infty\) means that the
depth indices are extended over all nonnegative integers.  Direct summation
of the geometric series gives
\begin{equation}\label{B:eq:generating}
\begin{aligned}
 D_\infty&=2(A^{-2}-\Delta^{-2}),\\
 B_\infty&=2(A^{-2}-(A\Delta)^{-1}),\\
 Q_{0,\infty}&=D_\infty-\frac{2(\Delta-A)}{A\Delta C},\\
 L_\infty&=2\left(\frac2{A^2}+\frac1{C^2}
                       -\frac1{AC}-\frac2{\Delta^2}\right).
\end{aligned}
\end{equation}
Indeed, the first two formulas follow from
$\sum_{p\ge0}(p+1)t^p=(1-t)^{-2}$ and the common-shift identity for the
minimum kernel.  For the pure cross block, setting $q=p+h$, $s=r+h$,
$h\ge1$, gives
\[
 \sum_{p,r\ge0}\min(p+1,r+1)(\xi\eta)^p(\zeta\omega)^r
     =\frac1{C(1-\Pi)},
\]
and
\[
 \sum_{h\ge1}(\eta^h-\xi^h)(\omega^h-\zeta^h)
     =\frac{(1-\Pi)(\Delta-A)}{A\Delta}.
\]
Realignment interchanges $A$ and $C$ and leaves $\Delta$ unchanged,
which yields the last line of \eqref{B:eq:generating}.

Write $\mathcal E_{\rm mix}:=M-B$.  By \Cref{B:lem:means}, the repeated
generating expression $\mathcal E_{\rm mix}(\xi,\xi,\zeta,\zeta)$
contains only total degrees at least $m$ in each pair of variables.
Moreover, positivity and \eqref{eq:sectors} imply the global bound
\begin{equation}\label{B:eq:global}
 0\le P_\pm[\xi,\eta]
 \le\frac8{(1-|\xi|^2)^2(1-|\eta|^2)^2},
 \qquad |\xi|,|\eta|<1.
\end{equation}
On conjugate nodes, the same-side four-form vanishes and
\begin{equation}\label{B:eq:repeated}
 P_-[\xi,\bar\xi]
 =Q_{0,m}[\xi,\bar\xi]
  +\mathcal E_{\rm mix}(\xi,\xi,\bar\xi,\bar\xi).
\end{equation}
If $\xi=\sqrt\rho e^{i\vartheta}$, the corresponding infinite baseline is
\begin{equation}\label{B:eq:repeated-baseline}
 Q_{0,\infty}[\xi,\bar\xi]
 =\frac{8\rho\sin^2\vartheta}
 {(1-\rho)^2((1-\rho)^2+4\rho\sin^2\vartheta)^2}.
\end{equation}

\begin{lemma}[Uniform finite tails]\label{B:lem:quant-tails}
For $0<\epsilon\le1$, set
\begin{equation}\label{B:eq:tail-budget}
 \mathcal T(m,\epsilon)=2^{32}\epsilon^{-8}e^{-\epsilon m/4}.
\end{equation}
If all four argument moduli are at most $e^{-\epsilon/2}$, then
$|L_m-L_\infty|\le\mathcal T(m,\epsilon)$.  For every positive
semidefinite normal-form Gram representation,
\begin{equation}\label{B:eq:conjugate-tail}
 |P_-[\xi,\bar\xi]-Q_{0,\infty}[\xi,\bar\xi]|
 \le\mathcal T(m,\epsilon)
 \qquad(|\xi|\le e^{-\epsilon/2}).
\end{equation}
\end{lemma}

\begin{proof}
The diagonal of the realigned correction vanishes, so the positive
semidefinite matrices $M$ and $B$ have the same diagonal.  For
$W=(p+1)(q+1)(r+1)(s+1)$, Cauchy--Schwarz gives
\[
 |(\mathcal E_{\rm mix})_{pq,rs}|\le4\sqrt W\le4W.
\]
The coefficients of the baseline kernels are bounded by a constant
multiple of $W$.  Every discarded finite-depth term contains at least one
index of size at least $m$.  Splitting the radial decay in half and using
$1-e^{-\epsilon/4}\ge\epsilon/8$ therefore gives
\[
 |L_m-L_\infty|,\ |Q_{0,m}-Q_{0,\infty}|
 \le2^{30}\epsilon^{-8}e^{-\epsilon m/4}.
\]
For the correction in \eqref{B:eq:repeated}, \Cref{B:lem:means} leaves
only total degrees at least $m$ in both variable pairs; the same summation
gives
\[
 |\mathcal E_{\rm mix}(\xi,\xi,\bar\xi,\bar\xi)|
 \le2^{26}\epsilon^{-8}e^{-\epsilon m/2}.
\]
Combining the two estimates proves the lemma.
\end{proof}

For the explicit threshold set
\begin{equation}\label{B:eq:dyadic-scale}
 k\ge16,\qquad \epsilon=2^{-k},\qquad m=\epsilon^{-2}=2^{2k}.
\end{equation}
Then
\begin{equation}\label{B:eq:dyadic-tails}
 \mathcal T\le2^{32+8k-2^{k-2}},\qquad
 \mathcal T\le\epsilon^{-2},\qquad
 \epsilon^4\mathcal T\le\epsilon.
\end{equation}

\begin{lemma}[Quantitative continuation]\label{B:lem:quant-continuation}
Let $P_-=\Gamma^\top\Gamma$ be a real Gram factorization, define
$\mathbf h(\xi,\eta)=\Gamma f(\xi,\eta)$, and set
\begin{equation}\label{B:eq:hepsilon}
 \mathbf h_\epsilon(\alpha,\beta)
 =\mathbf h\bigl(ie^{-\epsilon(\alpha+i\beta)},
        -ie^{-\epsilon(\alpha-i\beta)}\bigr).
\end{equation}
For the parameters in \eqref{B:eq:dyadic-scale},
\begin{equation}\label{B:eq:quant-continuation}
 \epsilon^4\norm{\mathbf h_\epsilon(\alpha,\beta)}^2
 \le2^{20}\epsilon^{2^{-17}}
\end{equation}
throughout
\begin{equation}\label{B:eq:speed-region}
 1\le\Re\alpha\le2,\quad |\Im\alpha|\le256,\quad
 |\Re\beta|\le256,\quad |\Im\beta|\le\tfrac12.
\end{equation}
The bound is uniform over the dimension of the Gram factorization.
\end{lemma}

\begin{proof}
From \eqref{B:eq:global} and $1-e^{-\epsilon/2}\ge\epsilon/4$,
\[
 \norm{\mathbf h_\epsilon(\alpha,\beta)}^2
 \le2^{20}\epsilon^{-4}
\]
whenever $\Re\alpha-|\Im\beta|\ge1/4$.  On the real conjugate slice,
\eqref{B:eq:repeated-baseline}, \Cref{B:lem:quant-tails}, and
\eqref{B:eq:dyadic-tails} improve this to
\[
 \norm{\mathbf h_\epsilon(\alpha,\beta)}^2
 \le2^{20}\epsilon^{-2}
\]
for $\alpha\in[1/2,1024+1/2]$ and $\beta\in[-512,512]$.

We apply the two-constant principle twice.  On the upper
$\alpha$-rectangle
$1/2\le\Re\alpha\le1024+1/2$, $0\le\Im\alpha\le512$, use
\[
 h_1(\alpha)=\sin\frac{\pi(\Re\alpha-1/2)}{1024}
 \frac{\sinh(\pi(512-\Im\alpha)/1024)}{\sinh(\pi/2)}.
\]
On $1\le\Re\alpha\le2$ and $0\le\Im\alpha\le256$ one has
$h_1>2^{-12}>2^{-13}$.  Repeating in the lower rectangle yields
\[
 \norm{\mathbf h_\epsilon(\alpha,\beta)}^2
 \le2^{20}\epsilon^{-4+2\cdot2^{-13}}
\]
for $1\le\Re\alpha\le2$, $|\Im\alpha|\le256$, and real
$\beta\in[-512,512]$.

For fixed such $\alpha$, apply the same argument in the upper and lower
$\beta$-rectangles with
\[
 h_2(\beta)=\cos\frac{\pi\Re\beta}{1024}
 \frac{\sinh(\pi(3/4-\Im\beta)/1024)}{\sinh(3\pi/4096)}.
\]
On $|\Re\beta|\le256$ and $|\Im\beta|\le1/2$, its relevant harmonic
measure is larger than $2^{-5}$.  Since
$2\cdot2^{-13}\cdot2^{-5}=2^{-17}$, the second comparison gives
\eqref{B:eq:quant-continuation}.
\end{proof}

